\section[Contrôles de la réalisation physico-dimensionnelle]{Domaine et obstructions des réalisations physique et dimensionnelle}
\label{sec:int68-certificado-historico}

La réalisation physico-dimensionnelle est soumise à dix familles de contrôle :
\begin{enumerate}
 \item concentration cylindrique, masse un et moments finis ;
 \item spectres \((9,5)\) du bloc APP et \((7,3)\) du bloc de mémoire ;
 \item identité \(\tanh^2(2\log\varphi)=5/9\) ;
 \item antiunitarité dans le modèle fini déclaré ;
 \item loi du produit semi-direct de Poincaré ;
 \item identités constitutives du vide ;
 \item norme scindée des modes nuls et clôture matérielle ;
 \item sélecteur \(\theta_*=(54/\pi)^2\) et point fixe autodual ;
 \item valeurs CGLMP de contrôle ;
 \item invariants \((396,729,43,387)\) de la filtration nonadique et
       domaine \((81,56,13,324)\), avec secteur nul séparé, de la Loi
       générale des masses.
\end{enumerate}

\begin{controlbox}[title={Portée exacte des contrôles}]
Les dix familles vérifient des identités, des domaines finis et des exemples
reproductibles. Elles ne remplacent pas la fonctorialité complète des quatre
composantes, ne construisent pas à elles seules le domaine commun, ne démontrent pas le
prolongement lisse de Cartan et n'élèvent pas les interfaces physiques ouvertes au rang de
théorèmes internes. La réalisation discrète conserve sa démonstration principale
dans la \cref{sec:int68-realizacion-fisica-discreta}.
\end{controlbox}

\subsection{Dépendances fonctorielles des réalisations spectrale, solénoïdale, de calibre, cohomologique et déterminantale}

Les contrôles précédents ne s'appliquent qu'après la construction d'APP, de TRIT, de TPK,
de l'état enrichi commun, de la double projection et de la droite dimensionnelle. La
concentration cylindrique et le spectre fini prouvent des propriétés de la
réalisation ; ils ne sélectionnent pas rétrospectivement l'état. La périodicité de
108 phases, l'atlas (81/56/13/324), l'involution constitutive et les
programmes cosmologiques conservent des domaines et des réfutateurs séparés.

Les identités de Dirac, de von Neumann, de CGLMP, de Cartan et de Poincaré sont utilisées
comme contrôles typés de l'opérateur d'entrelacement correspondant. Une coïncidence
numérique sans l'application de domaine, le codomaine et l'action sur les générateurs ne
constitue pas une réalisation physique.
